% Soluciones BE1 - Fibras Ópticas
\documentclass[11pt,a4paper]{article}

% Paquetes
\usepackage[utf8]{inputenc}
\usepackage[spanish]{babel}
\usepackage{amsmath}
\usepackage{amsfonts}
\usepackage{amssymb}
\usepackage{graphicx}
\usepackage{geometry}
\usepackage{fancyhdr}
\usepackage{hyperref}
\usepackage{enumitem}
\usepackage{siunitx}
\usepackage{xcolor}

\geometry{a4paper, margin=2.5cm}

% Configuración de encabezados
\pagestyle{fancy}
\fancyhf{}
\rhead{BE1 - Fibras Ópticas}
\lhead{IMT Atlantique}
\cfoot{\thepage}

\title{\textbf{BE1 : Fibras Ópticas}\\
\large Soluciones Detalladas\\
UE Tecnologías Ópticas - IMT Atlantique}
\author{Año 2025-2026}
\date{Enero 2026}

\begin{document}

\maketitle
\tableofcontents
\newpage

%=============================================================================
\section{Ejercicio 1: Modelo Geométrico de la Propagación}
%=============================================================================

\subsection{Datos del Problema}

Tenemos un guía óptico (fibra) con:
\begin{itemize}
    \item Índice de refracción del núcleo: $n_c$
    \item Índice de refracción de la cubierta: $n_g$
    \item Ángulo de inyección del rayo: $\theta$
    \item Longitud de la fibra: $L$
    \item Medio exterior: aire con $n = 1$
\end{itemize}

\subsection{Pregunta a) Condición para Propagación Guiada}

\textbf{Objetivo:} Encontrar la condición sobre $\theta$ que permite la propagación guiada.

\textbf{Análisis físico:}

Para que un rayo se propague por la fibra, debe sufrir reflexión total interna en la interfaz núcleo-cubierta. Esto requiere que el ángulo de incidencia en la interfaz sea mayor que el ángulo crítico.

\textbf{Paso 1: Aplicar la Ley de Snell en la entrada}

Cuando el rayo entra en la fibra desde el aire:
\begin{equation}
    n_{\text{aire}} \sin\theta = n_c \sin\theta'
\end{equation}

Como el aire tiene $n = 1$:
\begin{equation}
    \sin\theta = n_c \sin\theta'
\end{equation}

donde $\theta'$ es el ángulo de refracción dentro del núcleo respecto a la normal de la entrada.

\textbf{Paso 2: Relación geométrica}

El rayo refractado viaja en el núcleo con un ángulo $\alpha$ respecto al eje de la fibra. Por geometría:
\begin{equation}
    \alpha = 90^{\circ} - \theta'
\end{equation}

Por lo tanto:
\begin{equation}
    \sin\theta' = \cos\alpha
\end{equation}

\textbf{Paso 3: Condición de reflexión total interna}

Para que haya reflexión total en la interfaz núcleo-cubierta, el ángulo $\alpha$ debe ser mayor que el ángulo crítico $\alpha_c$:
\begin{equation}
    \sin\alpha_c = \frac{n_g}{n_c}
\end{equation}

\textbf{Paso 4: Combinar las condiciones}

Sustituyendo en la ecuación de Snell:
\begin{equation}
    \sin\theta = n_c \sin\theta' = n_c \cos\alpha
\end{equation}

Para el ángulo crítico:
\begin{equation}
    \sin\theta_{\max} = n_c \cos\alpha_c
\end{equation}

Usando la identidad trigonométrica $\cos\alpha_c = \sqrt{1 - \sin^2\alpha_c}$:
\begin{equation}
    \sin\theta_{\max} = n_c \sqrt{1 - \sin^2\alpha_c} = n_c \sqrt{1 - \frac{n_g^2}{n_c^2}}
\end{equation}

Simplificando:
\begin{equation}
    \sin\theta_{\max} = \sqrt{n_c^2 - n_g^2}
\end{equation}

\textbf{Condición para propagación guiada:}

El ángulo de inyección debe satisfacer:
\begin{equation}
    \sin\theta \leq \sqrt{n_c^2 - n_g^2}
\end{equation}

o equivalentemente:
\begin{equation}
    \theta \leq \arcsin\left(\sqrt{n_c^2 - n_g^2}\right)
\end{equation}

Este valor máximo se conoce como la \textbf{apertura numérica (NA)} de la fibra:
\begin{equation}
    \text{NA} = \sqrt{n_c^2 - n_g^2}
\end{equation}

\subsection{Pregunta b) Diferencia de Tiempos de Propagación}

\textbf{Objetivo:} Expresar $\Delta t$, la diferencia de tiempos de propagación máxima entre modos.

\textbf{Definición del índice de refracción relativo:}

Se define como:
\begin{equation}
    \Delta = \frac{n_c - n_g}{n_c}
\end{equation}

\textbf{Paso 1: Tiempo de propagación del rayo axial}

Un rayo que viaja paralelo al eje de la fibra recorre la distancia $L$ con velocidad $v = c/n_c$:
\begin{equation}
    t_{\min} = \frac{L}{c/n_c} = \frac{Ln_c}{c}
\end{equation}

\textbf{Paso 2: Tiempo de propagación del rayo más inclinado}

Un rayo que entra con el ángulo máximo $\theta_{\max}$ viaja en zigzag. La distancia total recorrida es mayor que $L$.

El rayo viaja con ángulo $\alpha_c$ respecto al eje. La distancia recorrida por reflexión es:
\begin{equation}
    d = \frac{L}{\cos\alpha_c}
\end{equation}

El tiempo correspondiente es:
\begin{equation}
    t_{\max} = \frac{d \cdot n_c}{c} = \frac{Ln_c}{c\cos\alpha_c}
\end{equation}

Como $\cos\alpha_c = \sin\theta' = n_g/n_c$:
\begin{equation}
    t_{\max} = \frac{Ln_c}{c} \cdot \frac{n_c}{n_g} = \frac{Ln_c^2}{cn_g}
\end{equation}

\textbf{Paso 3: Diferencia de tiempos}

\begin{align}
    \Delta t &= t_{\max} - t_{\min} \\
    &= \frac{Ln_c^2}{cn_g} - \frac{Ln_c}{c} \\
    &= \frac{Ln_c}{c}\left(\frac{n_c}{n_g} - 1\right) \\
    &= \frac{Ln_c}{c} \cdot \frac{n_c - n_g}{n_g}
\end{align}

Como $\Delta = (n_c - n_g)/n_c$, tenemos $n_c - n_g = \Delta \cdot n_c$:
\begin{equation}
    \Delta t = \frac{Ln_c}{c} \cdot \frac{\Delta n_c}{n_g}
\end{equation}

Para fibras con $n_g \approx n_c$ (aproximación débilmente guiante):
\begin{equation}
    \Delta t \approx \frac{Ln_c\Delta}{c}
\end{equation}

\textbf{Diferencia de tiempos de propagación:}
\begin{equation}
    \Delta t = \frac{Ln_c\Delta}{c}
\end{equation}

\subsection{Pregunta c) Fenómeno de Propagación}

\textbf{Respuesta:} Esta diferencia de tiempos hace referencia a la \textbf{dispersión modal} o \textbf{dispersión intermodal}.

\textbf{Explicación:}

En una fibra multimodo, diferentes rayos (o modos) pueden propagarse simultáneamente. Cada modo tiene un camino óptico diferente:
\begin{itemize}
    \item Los rayos axiales (modo fundamental) viajan la distancia más corta
    \item Los rayos más inclinados (modos de orden superior) recorren trayectorias más largas en zigzag
\end{itemize}

Esta diferencia de caminos causa que un pulso de luz se ensanche temporalmente durante la propagación, limitando la capacidad de transmisión de datos.

\subsection{Pregunta d) Débito Máximo}

\textbf{Objetivo:} Determinar el débito máximo considerando la dispersión modal.

\textbf{Razonamiento:}

Para que dos pulsos sucesivos no se superpongan en la recepción, el intervalo entre pulsos debe ser mayor que la dispersión temporal $\Delta t$.

\textbf{Paso 1: Condición de no solapamiento}

El tiempo mínimo entre bits es:
\begin{equation}
    T_{\min} = \Delta t
\end{equation}

\textbf{Paso 2: Débito máximo}

El débito (tasa de bits) es el inverso del tiempo entre bits:
\begin{equation}
    B_{\max} = \frac{1}{\Delta t}
\end{equation}

Sustituyendo la expresión de $\Delta t$:
\begin{equation}
    B_{\max} = \frac{c}{Ln_c\Delta}
\end{equation}

\textbf{Débito máximo:}
\begin{equation}
    B_{\max} = \frac{c}{Ln_c\Delta}
\end{equation}

Esta expresión muestra que:
\begin{itemize}
    \item El débito disminuye linealmente con la longitud $L$
    \item Fibras con menor $\Delta$ (diferencia de índices pequeña) permiten mayor débito
    \item Para aumentar el débito, se pueden usar fibras monomodo (eliminando la dispersión modal)
\end{itemize}

\subsection{Pregunta e) Cálculo Numérico}

\textbf{Datos:}
\begin{itemize}
    \item $L = 10$ km $= 10^4$ m
    \item $\Delta = 2 \times 10^{-2}$
    \item $n_c = 1.5$
    \item $c = 3 \times 10^8$ m/s
\end{itemize}

\textbf{Cálculo de $\Delta t$:}

Aplicamos la fórmula:
\begin{align}
    \Delta t &= \frac{Ln_c\Delta}{c} \\
    &= \frac{10^4 \times 1.5 \times 2 \times 10^{-2}}{3 \times 10^8} \\
    &= \frac{10^4 \times 1.5 \times 0.02}{3 \times 10^8} \\
    &= \frac{300}{3 \times 10^8} \\
    &= \frac{100}{10^8} \\
    &= 10^{-6} \text{ s}
\end{align}

\textbf{Resultado:}
\begin{equation}
    \Delta t = 1 \text{ }\mu\text{s} = 1 \text{ microsegundo}
\end{equation}

\textbf{Débito máximo correspondiente:}
\begin{equation}
    B_{\max} = \frac{1}{\Delta t} = \frac{1}{10^{-6}} = 10^6 \text{ bits/s} = 1 \text{ Mbit/s}
\end{equation}

\textbf{Interpretación:}

Para una fibra multimodo de 10 km con estos parámetros, la dispersión modal limita el débito a aproximadamente 1 Mbit/s. Este es un valor relativamente bajo para las telecomunicaciones modernas, lo que explica por qué se prefieren las fibras monomodo para transmisiones de larga distancia y alta velocidad.

%=============================================================================
\section{Ejercicio 2: Onda Electromagnética Monocromática}
%=============================================================================

\subsection{Datos del Problema}

\begin{itemize}
    \item Índice de refracción del núcleo: $n_c = 1.4555$
    \item Índice de refracción de la cubierta: $n_g = 1.4444$
    \item Longitud de onda: $\lambda = 1.55$ $\mu$m
\end{itemize}

\subsection{Condición para Fibra Monomodo}

\textbf{Objetivo:} Calcular el diámetro máximo del núcleo para que la fibra sea monomodo.

\textbf{Paso 1: Recordatorio teórico}

Una fibra óptica es monomodo cuando solo permite la propagación del modo fundamental LP$_{01}$. Esto ocurre cuando la frecuencia normalizada $V$ es menor que 2.405:
\begin{equation}
    V < 2.405
\end{equation}

\textbf{Paso 2: Definición de la frecuencia normalizada}

La frecuencia normalizada se define como:
\begin{equation}
    V = \frac{2\pi a}{\lambda}\text{NA}
\end{equation}

donde:
\begin{itemize}
    \item $a$ es el radio del núcleo
    \item $\lambda$ es la longitud de onda
    \item NA es la apertura numérica
\end{itemize}

\textbf{Paso 3: Cálculo de la apertura numérica}

\begin{align}
    \text{NA} &= \sqrt{n_c^2 - n_g^2} \\
    &= \sqrt{(1.4555)^2 - (1.4444)^2} \\
    &= \sqrt{2.1185 - 2.0863} \\
    &= \sqrt{0.0322} \\
    &= 0.1795
\end{align}

\textbf{Paso 4: Condición monomodo}

Para que la fibra sea monomodo:
\begin{equation}
    V \leq 2.405
\end{equation}

Sustituyendo la expresión de $V$:
\begin{equation}
    \frac{2\pi a}{\lambda}\text{NA} \leq 2.405
\end{equation}

\textbf{Paso 5: Despeje del radio máximo}

\begin{equation}
    a_{\max} = \frac{2.405 \lambda}{2\pi \text{NA}}
\end{equation}

\textbf{Paso 6: Cálculo numérico}

\begin{align}
    a_{\max} &= \frac{2.405 \times 1.55 \times 10^{-6}}{2\pi \times 0.1795} \\
    &= \frac{3.728 \times 10^{-6}}{1.128} \\
    &= 3.304 \times 10^{-6} \text{ m} \\
    &= 3.304 \text{ }\mu\text{m}
\end{align}

\textbf{Paso 7: Diámetro máximo}

El diámetro es el doble del radio:
\begin{align}
    d_{\max} &= 2a_{\max} \\
    &= 2 \times 3.304 \\
    &= 6.608 \text{ }\mu\text{m}
\end{align}

\textbf{Resultado:}

El diámetro máximo del núcleo para que la fibra sea monomodo a $\lambda = 1.55$ $\mu$m es:
\begin{equation}
    d_{\max} \approx 6.6 \text{ }\mu\text{m}
\end{equation}

\textbf{Observación:}

En la práctica, las fibras monomodo estándar G.652 tienen un diámetro de núcleo de aproximadamente 9 $\mu$m, diseñadas para operar en la ventana de 1310 nm y 1550 nm. El valor calculado es consistente con este diseño, considerando que queremos asegurar operación monomodo con un margen de seguridad.

%=============================================================================
\section{Ejercicio 3: Paquete de Información, Atenuación}
%=============================================================================

\subsection{Datos del Problema}

\begin{itemize}
    \item Potencia acoplada a la fibra: $P_{\text{in}} = 100$ $\mu$W
    \item Longitud de onda: $\lambda = 1550$ nm
    \item Atenuación medida: $A_{\text{total}} = 30$ dB
    \item Atenuación por kilómetro: $\alpha = 0.2$ dB/km
    \item Tipo de fibra: G.652 (monomodo estándar)
\end{itemize}

\subsection{Pregunta a) Potencia Medida en el Detector}

\textbf{Objetivo:} Calcular la potencia de salida en vatios (W) y en dBm.

\textbf{Paso 1: Recordatorio de la atenuación}

La atenuación en dB se define como:
\begin{equation}
    A[\text{dB}] = 10\log_{10}\left(\frac{P_{\text{in}}}{P_{\text{out}}}\right)
\end{equation}

\textbf{Paso 2: Despejar la potencia de salida}

De la ecuación anterior:
\begin{equation}
    \frac{P_{\text{in}}}{P_{\text{out}}} = 10^{A/10}
\end{equation}

Por lo tanto:
\begin{equation}
    P_{\text{out}} = \frac{P_{\text{in}}}{10^{A/10}}
\end{equation}

\textbf{Paso 3: Conversión de la potencia de entrada}

Primero convertimos la potencia de entrada:
\begin{align}
    P_{\text{in}} &= 100 \text{ }\mu\text{W} \\
    &= 100 \times 10^{-6} \text{ W} \\
    &= 10^{-4} \text{ W}
\end{align}

\textbf{Paso 4: Cálculo de la potencia de salida en W}

\begin{align}
    P_{\text{out}} &= \frac{10^{-4}}{10^{30/10}} \\
    &= \frac{10^{-4}}{10^3} \\
    &= 10^{-7} \text{ W} \\
    &= 0.1 \text{ }\mu\text{W} \\
    &= 100 \text{ nW}
\end{align}

\textbf{Potencia de salida en vatios:}
\begin{equation}
    P_{\text{out}} = 10^{-7} \text{ W} = 100 \text{ nW}
\end{equation}

\textbf{Paso 5: Conversión a dBm}

La potencia en dBm se calcula con:
\begin{equation}
    P[\text{dBm}] = 10\log_{10}\left(\frac{P[\text{W}]}{10^{-3}}\right)
\end{equation}

Para la potencia de entrada:
\begin{align}
    P_{\text{in}}[\text{dBm}] &= 10\log_{10}\left(\frac{10^{-4}}{10^{-3}}\right) \\
    &= 10\log_{10}(10^{-1}) \\
    &= -10 \text{ dBm}
\end{align}

Para la potencia de salida:
\begin{align}
    P_{\text{out}}[\text{dBm}] &= P_{\text{in}}[\text{dBm}] - A[\text{dB}] \\
    &= -10 - 30 \\
    &= -40 \text{ dBm}
\end{align}

Verificación:
\begin{align}
    P_{\text{out}}[\text{dBm}] &= 10\log_{10}\left(\frac{10^{-7}}{10^{-3}}\right) \\
    &= 10\log_{10}(10^{-4}) \\
    &= -40 \text{ dBm}
\end{align}

\textbf{Potencia de salida en dBm:}
\begin{equation}
    P_{\text{out}} = -40 \text{ dBm}
\end{equation}

\subsection{Pregunta b) Longitud Estimada de la Fibra}

\textbf{Objetivo:} Calcular la longitud de la fibra a partir de la atenuación.

\textbf{Paso 1: Relación entre atenuación total y longitud}

La atenuación total es el producto de la atenuación por unidad de longitud por la longitud:
\begin{equation}
    A_{\text{total}} = \alpha \cdot L
\end{equation}

\textbf{Paso 2: Despeje de la longitud}

\begin{equation}
    L = \frac{A_{\text{total}}}{\alpha}
\end{equation}

\textbf{Paso 3: Cálculo numérico}

\begin{align}
    L &= \frac{30 \text{ dB}}{0.2 \text{ dB/km}} \\
    &= 150 \text{ km}
\end{align}

\textbf{Longitud estimada de la fibra:}
\begin{equation}
    L = 150 \text{ km}
\end{equation}

\textbf{Interpretación:}

Esta longitud de 150 km es típica para enlaces de telecomunicaciones de media distancia. A 1550 nm, la fibra G.652 presenta su mínima atenuación (alrededor de 0.2 dB/km), lo que permite enlaces largos sin necesidad de amplificación. Para distancias mayores, se requerirían amplificadores ópticos (EDFA) o repetidores.

%=============================================================================
\section{Ejercicio 4: Paquete de Información, Dispersión Cromática}
%=============================================================================

Este ejercicio requiere implementación de código para simular la propagación en fibra óptica. A continuación se presenta la solución conceptual y las ecuaciones necesarias.

\subsection{Introducción}

El objetivo es simular la propagación de señales en fibra óptica considerando la dispersión cromática, utilizando la ecuación no lineal de Schrödinger (NLSE) en su forma simplificada (solo dispersión).

\subsection{Pregunta a) Generación de Secuencia Binaria}

\textbf{Objetivo:} Crear una secuencia binaria pseudo-aleatoria.

\textbf{Procedimiento:}

En cualquier lenguaje de programación (Python, MATLAB, etc.), generar una secuencia aleatoria de $n_{\text{bits}}$ bits (0 y 1).

Ejemplo conceptual:
\begin{itemize}
    \item Definir el número de bits: $n_{\text{bits}} = 100$
    \item Generar valores aleatorios entre 0 y 1
    \item Asignar 0 si el valor es menor que 0.5, y 1 en caso contrario
\end{itemize}

\subsection{Pregunta b) Modulación NRZ}

\textbf{Objetivo:} Convertir la secuencia binaria en una señal temporal continua.

\textbf{Parámetros:}
\begin{itemize}
    \item Tasa de muestreo: $F_s$ (Hz)
    \item Débito: $B$ (bits/s)
    \item Amplitud: $+1$ V para bit 0, $+2$ V para bit 1
\end{itemize}

\textbf{Procedimiento:}

La duración de cada bit es:
\begin{equation}
    T_{\text{bit}} = \frac{1}{B}
\end{equation}

El número de muestras por bit es:
\begin{equation}
    N_{\text{muestras}} = \frac{F_s}{B}
\end{equation}

Para modulación NRZ (Non-Return-to-Zero):
\begin{itemize}
    \item Cada bit mantiene su amplitud durante toda su duración
    \item Bit 0 $\rightarrow$ amplitud = 1 V
    \item Bit 1 $\rightarrow$ amplitud = 2 V
\end{itemize}

El término DC (componente continua) se añade para emular un láser modulado directamente.

\subsection{Pregunta c) Propagación en Fibra}

\textbf{Ecuación diferencial:}

La envolvente de la señal evoluciona según:
\begin{equation}
    \frac{\partial A(z,t)}{\partial z} = -\frac{i\beta_2}{2}\frac{\partial^2 A(z,t)}{\partial t^2}
\end{equation}

donde $\beta_2$ es el parámetro de dispersión de segundo orden.

\textbf{Relación con el coeficiente de dispersión:}

\begin{equation}
    \beta_2 = -\frac{D\lambda^2}{2\pi c}
\end{equation}

Para fibra G.652 a $\lambda = 1550$ nm:
\begin{itemize}
    \item $D \approx 17$ ps/(nm·km)
    \item Longitud de onda: $\lambda = 1550$ nm
    \item Velocidad de la luz: $c = 3 \times 10^8$ m/s
\end{itemize}

\textbf{Cálculo de $\beta_2$:}

\begin{align}
    \beta_2 &= -\frac{17 \times 10^{-12} \times (1550 \times 10^{-9})^2}{2\pi \times 3 \times 10^8 \times 10^3} \\
    &= -\frac{17 \times 10^{-12} \times 2.4025 \times 10^{-12}}{1.885 \times 10^{12}} \\
    &\approx -21.7 \times 10^{-27} \text{ s}^2\text{/m}
\end{align}

\textbf{Método de resolución - Split-Step Fourier:}

La solución analítica en el dominio de Fourier es:
\begin{equation}
    \tilde{A}(z,\omega) = \tilde{A}(0,\omega) \exp\left(\frac{i\beta_2\omega^2z}{2}\right)
\end{equation}

Algoritmo:
\begin{enumerate}
    \item Transformada de Fourier de la señal de entrada: $A(0,t) \rightarrow \tilde{A}(0,\omega)$
    \item Multiplicar por la función de transferencia: $H(\omega) = \exp(i\beta_2\omega^2 L/2)$
    \item Transformada inversa de Fourier para obtener: $A(L,t)$
\end{enumerate}

\subsection{Pregunta d) Detección Fotónica}

\textbf{Modelo simplificado:}

El fotodetector genera una corriente proporcional a la intensidad óptica:
\begin{equation}
    I_{\text{foto}}(t) \propto |A(L,t)|^2
\end{equation}

El bloque DC elimina la componente continua:
\begin{equation}
    I_{\text{AC}}(t) = I_{\text{foto}}(t) - \langle I_{\text{foto}}(t) \rangle
\end{equation}

\subsection{Pregunta e) Casos de Uso}

\textbf{Parámetros de simulación:}
\begin{itemize}
    \item $\lambda = 1550$ nm
    \item Fibra G.652: $D = 17$ ps/(nm·km)
    \item Tasa de muestreo: 32 muestras por símbolo
\end{itemize}

\textbf{Análisis cualitativo:}

La dispersión cromática causa ensanchamiento temporal del pulso. La longitud de dispersión característica es:
\begin{equation}
    L_D = \frac{T_0^2}{|\beta_2|}
\end{equation}

donde $T_0$ es la duración del pulso.

Para un débito $B$, la duración del bit es $T_{\text{bit}} = 1/B$.

\textbf{Estimación de longitud máxima:}

Para evitar interferencia entre símbolos (ISI), se requiere que el ensanchamiento sea menor que el periodo del bit:
\begin{equation}
    \Delta T < T_{\text{bit}}
\end{equation}

El ensanchamiento por dispersión es aproximadamente:
\begin{equation}
    \Delta T \approx D \cdot \Delta\lambda \cdot L
\end{equation}

Para una fuente con ancho espectral $\Delta\lambda$:

\textbf{Débito 1 Gb/s:}
\begin{itemize}
    \item $T_{\text{bit}} = 1$ ns
    \item Longitud máxima estimada: $L_{\max} \approx 100$ km
\end{itemize}

\textbf{Débito 10 Gb/s:}
\begin{itemize}
    \item $T_{\text{bit}} = 100$ ps
    \item Longitud máxima estimada: $L_{\max} \approx 10$ km
\end{itemize}

\textbf{Débito 50 Gb/s:}
\begin{itemize}
    \item $T_{\text{bit}} = 20$ ps
    \item Longitud máxima estimada: $L_{\max} \approx 2$ km
\end{itemize}

\textbf{Débito 200 Gb/s:}
\begin{itemize}
    \item $T_{\text{bit}} = 5$ ps
    \item Longitud máxima estimada: $L_{\max} \approx 500$ m
\end{itemize}

\subsection{Pregunta f) Régimen de Dispersión Normal}

En el régimen de dispersión normal (por ejemplo, a $\lambda = 1310$ nm para fibra G.652):
\begin{itemize}
    \item El signo de $\beta_2$ cambia
    \item El ensanchamiento del pulso sigue siendo proporcional a la distancia
    \item La forma del ensanchamiento es diferente (chirp positivo vs negativo)
\end{itemize}

El comportamiento general es similar, pero la fase acumulada es opuesta.

\subsection{Pregunta g) Técnicas de Compensación}

\textbf{Compensación de dispersión cromática:}

\begin{enumerate}
    \item \textbf{Fibra de compensación de dispersión (DCF):}
    \begin{itemize}
        \item Usar fibras con dispersión negativa para compensar la acumulada
        \item Ventaja: totalmente óptico
        \item Desventaja: introduce pérdidas adicionales
    \end{itemize}
    
    \item \textbf{Redes de Bragg (FBG):}
    \begin{itemize}
        \item Filtros que introducen retardos dependientes de la frecuencia
        \item Compactos y eficientes
    \end{itemize}
    
    \item \textbf{Ecualización electrónica (DSP):}
    \begin{itemize}
        \item Procesamiento digital de señal en el receptor
        \item Muy flexible, permite compensación adaptativa
        \item Usado en sistemas coherentes modernos
    \end{itemize}
    
    \item \textbf{Pre-compensación (chirp del láser):}
    \begin{itemize}
        \item Modular el láser para introducir chirp opuesto
        \item Compensación en la fuente
    \end{itemize}
    
    \item \textbf{Multiplexación por división de longitud de onda (WDM):}
    \begin{itemize}
        \item Usar múltiples canales con menor débito cada uno
        \item Reduce el impacto de la dispersión por canal
    \end{itemize}
\end{enumerate}

\subsection{Pregunta h) Métodos de Simulación Avanzados}

\textbf{Ecuación NLSE completa:}

Para incluir auto-modulación de fase (SPM):
\begin{equation}
    \frac{\partial A}{\partial z} = -\frac{\alpha}{2}A - \frac{i\beta_2}{2}\frac{\partial^2 A}{\partial t^2} + i\gamma|A|^2A
\end{equation}

donde:
\begin{itemize}
    \item $\alpha$: atenuación
    \item $\beta_2$: dispersión
    \item $\gamma$: coeficiente no lineal
\end{itemize}

\textbf{Método Split-Step Fourier (SSF):}

Se alterna entre dominio temporal (no linealidad) y frecuencial (dispersión):
\begin{enumerate}
    \item Dividir la fibra en segmentos pequeños $\Delta z$
    \item En cada segmento:
    \begin{itemize}
        \item Aplicar operador no lineal: $A \rightarrow A \exp(i\gamma|A|^2\Delta z)$
        \item Aplicar operador dispersivo en dominio de Fourier
    \end{itemize}
    \item Iterar hasta completar la longitud total
\end{enumerate}

\textbf{Precisión del método:}

El tamaño del paso debe satisfacer:
\begin{equation}
    \Delta z < \min(L_D, L_{\text{NL}})
\end{equation}

donde $L_{\text{NL}} = 1/(\gamma P_0)$ es la longitud no lineal.

%=============================================================================
\section{Ejercicio 5: Caso de Uso - Cable Óptico Submarino}
%=============================================================================

\subsection{Análisis de la Estructura del Cable}

Un cable óptico submarino está diseñado para proteger las fibras ópticas en condiciones extremas: alta presión, agua salada, corrientes marinas, y posibles daños mecánicos.

\textbf{Estructura típica (de exterior a interior):}
\begin{enumerate}
    \item Polietileno (protección exterior)
    \item Cinta Mylar (aislamiento)
    \item Alambres de acero trenzados (resistencia mecánica)
    \item Barrera de aluminio contra agua
    \item Policarbonato (aislamiento adicional)
    \item Tubo de cobre o aluminio
    \item Gel de petróleo
    \item Fibras ópticas
\end{enumerate}

\subsection{Pregunta a) Rol de la Estructura n°6}

\textbf{Estructura n°6: Tubo de cobre o aluminio}

\textbf{Funciones principales:}

\begin{enumerate}
    \item \textbf{Conducción eléctrica:}
    \begin{itemize}
        \item Transporta energía eléctrica para alimentar los repetidores/amplificadores submarinos
        \item Tensión típica: varios kilovoltios DC
        \item Corriente: decenas de amperios
        \item Los repetidores están espaciados cada 50-100 km
    \end{itemize}
    
    \item \textbf{Protección mecánica adicional:}
    \begin{itemize}
        \item Proporciona rigidez estructural
        \item Protege las fibras de esfuerzos de flexión
    \end{itemize}
    
    \item \textbf{Barrera contra humedad:}
    \begin{itemize}
        \item Junto con las otras capas, previene la entrada de agua
        \item El cobre/aluminio no se corroe fácilmente bajo el polietileno
    \end{itemize}
\end{enumerate}

\textbf{Importancia en sistemas submarinos:}

En un cable submarino transatlántico de 6000 km, puede haber 60-120 repetidores que necesitan alimentación continua. El tubo conductor es esencial para este propósito.

\subsection{Pregunta b) Protección contra el Agua}

\textbf{¿Por qué proteger las fibras ópticas del agua?}

Aunque el vidrio (sílice) es químicamente estable, el agua causa varios problemas graves:

\begin{enumerate}
    \item \textbf{Degradación mecánica (fatiga estática):}
    \begin{itemize}
        \item Las moléculas de agua atacan los enlaces Si-O-Si en la superficie del vidrio
        \item Reacción química: Si-O-Si + H$_2$O $\rightarrow$ 2Si-OH
        \item Esto crea micro-fisuras que se propagan bajo tensión
        \item Resultado: ruptura de la fibra a largo plazo (años)
    \end{itemize}
    
    \item \textbf{Aumento de la atenuación:}
    \begin{itemize}
        \item El agua absorbe fuertemente en el infrarrojo cercano
        \item Pico de absorción del OH$^-$ a 1383 nm (en la ventana de telecomunicaciones)
        \item Incluso trazas de humedad aumentan significativamente las pérdidas
        \item La atenuación puede pasar de 0.2 dB/km a varios dB/km
    \end{itemize}
    
    \item \textbf{Degradación del recubrimiento:}
    \begin{itemize}
        \item El recubrimiento de acrilato puede hincharse o degradarse
        \item Pérdida de las propiedades mecánicas de protección
    \end{itemize}
    
    \item \textbf{Hidrógeno molecular:}
    \begin{itemize}
        \item En ambiente húmedo, puede formarse H$_2$
        \item El H$_2$ difunde en el vidrio y causa pérdidas adicionales
        \item Efecto particularmente problemático en la ventana de 1550 nm
    \end{itemize}
\end{enumerate}

\textbf{Soluciones implementadas:}

\begin{itemize}
    \item Múltiples barreras herméticas (aluminio, acero)
    \item Gel hidrófobo (gel de petróleo) que repele el agua
    \item Recubrimientos especiales resistentes a la humedad
    \item Hermetic sealing en los empalmes y conectores
\end{itemize}

\textbf{Consecuencias de la penetración de agua:}

En cables submarinos sin protección adecuada:
\begin{itemize}
    \item Vida útil reducida de 25 años a menos de 5 años
    \item Fallos catastróficos por ruptura de fibras
    \item Interrupción del servicio (muy costosa en cables submarinos)
\end{itemize}

\subsection{Pregunta c) Componentes Ópticos en Repetidores}

\textbf{Fenómenos de propagación a compensar:}

En transmisión por fibra monomodo de larga distancia, los principales fenómenos limitantes son:

\begin{enumerate}
    \item \textbf{Atenuación:}
    \begin{itemize}
        \item Pérdida de potencia óptica
        \item Típicamente 0.2 dB/km a 1550 nm
        \item En 100 km: 20 dB de pérdida (factor 100)
    \end{itemize}
    
    \item \textbf{Dispersión cromática:}
    \begin{itemize}
        \item Ensanchamiento temporal de los pulsos
        \item Limitante para altos débitos
    \end{itemize}
    
    \item \textbf{Dispersión por modo de polarización (PMD):}
    \begin{itemize}
        \item Diferencia de velocidad entre polarizaciones
        \item Importante en fibras antiguas
    \end{itemize}
    
    \item \textbf{Efectos no lineales:}
    \begin{itemize}
        \item Auto-modulación de fase (SPM)
        \item Mezclado de cuatro ondas (FWM)
        \item Dispersión Raman estimulada (SRS)
    \end{itemize}
\end{enumerate}

\textbf{Componentes necesarios en cada repetidor:}

\begin{enumerate}
    \item \textbf{Amplificador óptico (EDFA - Erbium Doped Fiber Amplifier):}
    \begin{itemize}
        \item Función: Compensar la atenuación acumulada
        \item Tecnología: Fibra dopada con erbio bombeada a 980 nm o 1480 nm
        \item Ganancia típica: 20-30 dB
        \item Ventaja: Amplifica todas las longitudes de onda WDM simultáneamente
        \item Banda de operación: 1530-1565 nm (banda C)
    \end{itemize}
    
    \item \textbf{Módulo de compensación de dispersión (DCM):}
    \begin{itemize}
        \item Función: Compensar la dispersión cromática acumulada
        \item Tecnologías posibles:
        \begin{itemize}
            \item Fibra de compensación de dispersión (DCF): fibra con $D$ negativo
            \item Redes de Bragg en fibra (FBG): filtros dispersivos
            \item Compensadores basados en cristales fotónicos
        \end{itemize}
        \item Se coloca típicamente después del amplificador
    \end{itemize}
    
    \item \textbf{Filtros ópticos pasa-banda:}
    \begin{itemize}
        \item Función: Eliminar el ruido ASE (Amplified Spontaneous Emission) fuera de banda
        \item El ASE se acumula en cascada de amplificadores
        \item Mejora la relación señal a ruido óptica (OSNR)
    \end{itemize}
    
    \item \textbf{Ecualizador de ganancia (GEQ):}
    \begin{itemize}
        \item Función: Igualar la potencia entre canales WDM
        \item La ganancia del EDFA no es perfectamente plana
        \item Evita acumulación de desequilibrio de potencia
    \end{itemize}
    
    \item \textbf{Monitor de desempeño óptico (OPM):}
    \begin{itemize}
        \item Función: Supervisión y diagnóstico
        \item Mide potencia por canal, OSNR
        \item Permite gestión remota del repetidor
    \end{itemize}
    
    \item \textbf{Aisladores ópticos:}
    \begin{itemize}
        \item Función: Prevenir reflexiones hacia atrás
        \item Protegen al amplificador de realimentación
        \item Basados en efecto Faraday (no recíproco)
    \end{itemize}
\end{enumerate}

\textbf{Configuración típica de un repetidor (en orden):}

\begin{equation}
    \text{Entrada} \rightarrow \text{Aislador} \rightarrow \text{EDFA} \rightarrow \text{DCM} \rightarrow \text{Filtro} \rightarrow \text{GEQ} \rightarrow \text{Aislador} \rightarrow \text{Salida}
\end{equation}

\textbf{Parámetros importantes:}

\begin{itemize}
    \item Espaciamiento entre repetidores: 50-100 km
    \item Número de canales WDM: 40-100 canales
    \item Débito por canal: 10-100 Gb/s
    \item Capacidad total: varios Tb/s
    \item Vida útil del repetidor: 25 años sin mantenimiento
\end{itemize}

\textbf{Consideraciones para sistemas submarinos:}

\begin{enumerate}
    \item \textbf{Fiabilidad extrema:}
    \begin{itemize}
        \item Los repetidores deben funcionar sin fallo durante 25 años
        \item Redundancia en componentes críticos
        \item Componentes de grado espacial
    \end{itemize}
    
    \item \textbf{Eficiencia energética:}
    \begin{itemize}
        \item La potencia se transmite por el cable desde tierra
        \item Consumo típico: 10-20 W por repetidor
        \item Bombeo eficiente de los EDFAs
    \end{itemize}
    
    \item \textbf{Gestión térmica:}
    \begin{itemize}
        \item Disipación de calor en ambiente de alta presión
        \item Temperatura estable para operación óptima
    \end{itemize}
    
    \item \textbf{Escalabilidad:}
    \begin{itemize}
        \item Diseño que permita futuros upgrades de capacidad
        \item Ancho de banda del amplificador > ancho de banda usado inicialmente
    \end{itemize}
\end{enumerate}

\textbf{Ejemplo: Cable transatlántico}

Un cable submarino de 6000 km entre Europa y América del Norte tendría:
\begin{itemize}
    \item 60-100 repetidores espaciados cada 60-100 km
    \item Capacidad: 10-20 Tb/s
    \item Costo total: 200-300 millones de euros
    \item Tiempo de instalación: varios meses
    \item Vida útil: 25 años
\end{itemize}

\textbf{Evolución tecnológica:}

Los cables submarinos modernos incorporan:
\begin{itemize}
    \item Amplificadores Raman distribuidos (además de EDFAs)
    \item Multiplexación por división de modo (MDM) experimental
    \item Fibras de área efectiva grande (para reducir no linealidades)
    \item Detección coherente en el receptor (en tierra)
    \item Compensación digital de dispersión (en tierra)
\end{itemize}

\end{document}
